% !TeX root = Tools-from-complex-geometry.tex
\section{Tools from complex geometry}

The main references of this section are \cite{DeMarco-2003-DynamicsRationalMaps,DeMarco-2018-DynamicalModuliSpaces,Dinh-Sibony-2010-DynamicsSeveralComplex}; for basic knowledge in complex geometry, I refer to \cite{Barlet-Magnusson-2019-ComplexAnalyticCycles,Demailly-2012-ComplexAnalyticDifferential}. \Yang{}


\begin{notation}[label=\labelNFM{notation:complex-geometry}, status=draft, tag=0000000]
Throughout this section, we work over $\bbC$.
Let us fix some notations we will frequently use later.
	\begin{itemize}
		\item On the field $\bbC$ of complex numbers, we always fix a subfield $\bbR\subset \bbC$ of real numbers, a standard \latexdeleted{basis}\latexmath{real vector space decomposition} $\bbC = \bbR \oplus \bbR \im$, $\upi = \im$, and the standard \latexdeleted{Eculidean}\latexgrammar{Euclidean} topology.
		\item A complex analytic scheme is \latexdeleted{a locally cut off by holomorphic function, of finite type,  ...}
		      \begin{latexmathchange}
			      a ringed space locally isomorphic to $(A,(\scrO_\Omega/\calI)|_A)$, where $\Omega\subset\bbC^N$ is open, $\calI\subset\scrO_\Omega$ is a coherent ideal sheaf, and $A$ is its zero set.
			      Thus it is locally defined by finitely many holomorphic equations, with nilpotents allowed in its structure sheaf.
		      \end{latexmathchange}
		      {\raggedright
		      A complex analytic space is a reduced \latexdeleted{...}\latexmath{complex analytic scheme,} and a complex \latexdeleted{anayltic}\latexgrammar{analytic} variety is \latexmath{a reduced,} irreducible\latexmath{ complex analytic scheme}.
		      \par}
		      \begin{latexmathchange}
			      These are our terminology conventions; see \cite[Chapter~II, \S 5 and \S 9]{Demailly-2012-ComplexAnalyticDifferential}.
		      \end{latexmathchange}
		      {\raggedright
		      A manifold is \latexdeleted{awalys}\latexgrammar{always} connected \latexdeleted{without specific}\latexgrammar{unless otherwise stated}.
		      And all \latexdeleted{space}\latexgrammar{spaces} (with analytic topology) are Hausdorff and \latexdeleted{the} second countable.
		      \latexdeleted{For even}\latexgrammar{Even for} algebraic varieties, we only consider \latexdeleted{its closed point}\latexgrammar{their closed points}.
		      \par}
		\item Let $X$ be a complex analytic space\latexdeleted{,}\latexgrammar{.}
		      \latexdeleted{its sheaf of holomorphic funciton}\latexgrammar{We use the notation} $\scrO_X$, $\cinf_{X,\bbK}$, $\scrC^p_{X,\bbK}$\latexdeleted{, means the sheaf...}
		      \begin{latexmathchange}
			      for the sheaves of holomorphic functions, smooth $\bbK$-valued functions, and $C^p$ $\bbK$-valued functions, respectively, where $\bbK\in\{\bbR,\bbC\}$ and $p\geq0$ is an integer.
			      On a singular space, smooth or $C^p$ means locally the restriction of an ambient smooth or $C^p$ function under a local closed embedding into some $\bbC^N$.
		      \end{latexmathchange}
		\item \latexdeleted{sheaf of forms and currents,}\latexgrammar{For the sheaves of forms and currents,} we will use the notation $\latexdeletedmath{\scrA^..}\latexaddedmath{\scrA_X^{p,q}}$ and $\latexdeletedmath{\scrA_..}\latexaddedmath{\scrA_{p,q}^X}$.
		      \latexdeleted{note this maybe not standard.}\latexgrammar{This notation may not be standard.}
		      \begin{latexmathchange}
			      Forms are indexed by bidegree, and currents by bidimension; the total degree or dimension is denoted by $\scrA_X^k$ or $\scrA_k^X$.
			      On a complex manifold of dimension $d$, a current of bidimension $(p,q)$ has bidegree $(d-p,d-q)$.
		      \end{latexmathchange}
		\item $\upd^c$ is normalized by \latexdeleted{...}
		      \[
			      \latexaddedmath{\upd^c=\frac{1}{2\pi\upi}(\partial-\partialbar),\qquad
			      \ddc=\frac{\upi}{\pi}\partial\partialbar,}
		      \]
		      \latexdeleted{as}\latexgrammar{as in} \cite{Demailly-2012-ComplexAnalyticDifferential}.
		      \Yang{}
	\end{itemize}

	% \toai{finish and polish this paragraph} -- resolved in place.
	\Yang{}
\end{notation}


\begin{theorem}[label=\labelNFM{thm:partition-of-unit}, status=draft, tag=0000000, title={Partition of \texorpdfstring{\latexdeleted{unit}\latexgrammar{unity}}{unity}}]
	Let $X$ be a complex analytic space and $\{U_i\}_{\latexaddedmath{i\in I}}$ be an open cover of $X$.
	\latexdeleted{Then, there exist a locally finite refine $\{V_i\}$ of $\{U_i\}$ and $\varphi_i \in \cinf_{X,\bbR}(X)$ with $\Supp \varphi_i \subset V_i$ and $\sum \varphi_i = 1$.}
	\begin{latexmathchange}
		Then there exist a locally finite open refinement $\{V_j\}_{j\in J}$, a map $\iota:J\to I$, and functions $\varphi_j\in\cinf_{X,\bbR}(X)$ such that
		\begin{align*}
			V_j&\subset U_{\iota(j)},\qquad \Supp\varphi_j\subset V_j,\qquad 0\leq\varphi_j\leq1,\\
			\sum_{j\in J}\varphi_j&=1.
		\end{align*}
		The sum is locally finite; the functions form a smooth partition of unity subordinate to the refinement.
		For manifolds, see \cite[Theorem~4.1.7]{Dinh-Sibony-2005-IntroductionTheoryCurrents}.
		For reduced analytic spaces, the same construction uses ambient smooth bump functions in local embeddings and paracompactness of $X$.
	\end{latexmathchange}
	% \toai{finish and polish it} -- resolved in place.
\end{theorem}

\Yang{}









\subsection[status=draft,tag=0000000,label=\labelNFM{subsec:forms-and-currents}]{Review of forms and currents}

Recall that on a complex mainfolds $X$, or more generally a complex analytic space, we have the sheaves $\scrA_X^k$ of $\scrC^\infty$-differential forms of degree $k$.
On a regular complex coordinates $(U;z_1,\cdots,z_d)$,
\[ \scrA^k_X(U) = \left\{ \sum_{|I|+|J| = k} f_{I,J} \upd z_I \wedge \upd \bar{z}_J \;\middle|\; f_{I,J} \in \cinf_{X,\bbC}(U) \right\} = \bigwedge^k \scrA^1_X(U). \]
There is a $\bbR$-linear map, induced by the conjugation on $\bbC$, on $\scrA_X^k$ given by
\[ \alpha = \sum f_{I,J} \upd z_I \wedge \upd \bar{z}_J \quad \mapsto \quad \bar{\alpha} \coloneqq \sum \overline{f_{I,J}} \upd \bar{z}_I \wedge \upd z_J. \]
Set the sheaves $\scrA_X^{p,q}$ of $\scrC^\infty$-differential forms of bidegree $(p,q)$ by
\[ \scrA_X^{p,q}(U) = \left\{ \sum_{|I|=p, |J|=q} f_{I,J} \upd z_I \wedge \upd \bar{z}_J \;\middle|\; f_{I,J} \in \cinf_{X,\bbC}(U) \right\}. \]
Then we have
\[ \scrA_X^k = \bigoplus_{p+q = k} \scrA_X^{p,q}, \quad \scrA_X^{0,1} = \overline{\scrA_{X}^{1,0}}, \quad \scrA_X^{p,q} = \bigwedge^p \scrA_X^{1,0} \otimes \bigwedge^q \scrA_X^{0,1}. \]
Denote $\scrA_X \coloneqq \bigoplus_{k} \scrA_X^k = \bigoplus_{k} \bigoplus_{p+q = k} \scrA_X^{p,q}$.
All these sheaves are $\cinf_{X,\bbC}$-modules.
And the usual differential map $\upd: \scrA^k_X \to \scrA_X^{k+1}$ decomposition as $\upd = \partial + \partialbar$ with
\begin{align*}
	\partial    & : \scrA^{p,q}_X \to \scrA^{p+1,q}_X, \quad \sum f_{I,J} \upd z_I \wedge \upd \bar{z}_J \mapsto \sum \partial_{z_i}(f_{I,J})  \upd z_i \wedge \upd z^I \wedge \upd \bar{z}^J,            \\
	\partialbar & : \scrA^{p,q}_X \to \scrA^{p,q+1}_X, \quad \sum f_{I,J} \upd z_I \wedge \upd \bar{z}_J \mapsto \sum \partial_{\bar{z_i}}(f_{I,J}) \upd z_I \wedge \upd \bar{z_i} \wedge \upd \bar{z}_J.
\end{align*}
By $\upd^2 = 0$, we get $\partial^2 = \partialbar^2 = 0$ and $\partial\partialbar+\partialbar\partial = 0$.
Set
\[
	\upd^c\coloneqq\frac{1}{2\pi \upi}(\partial-\partialbar): \scrA_X^{k} \mapsto \scrA_X^{k+1}.
\]
Under this normalization, we have
\[
	\ddc = -\frac{1}{\pi\upi}\partial\partialbar = \frac{\upi}{\pi}\partial\partialbar: \scrA_X^{p,q} \longrightarrow \scrA_X^{p+1,q+1}.
\]
For details, see \cite{Barlet-Magnusson-2019-ComplexAnalyticCycles,Demailly-2012-ComplexAnalyticDifferential}.

There is a natural way to endow topology on $\scrA_X^{p,q}(U)$ making it as a Fr\'echet space.

\begin{construction}[label=\labelNFM{constr:topo-on-forms}, status=draft, tag=0000000]
	Let $X$ be a complex manifold of complex dimension $d$, $U\subset X$ an open set, and $0\leq p,q\leq d$.
	A \emph{compact subset with a coordinate} means a triple $(V,z,K)$, where $(V,z)$ is a holomorphic coordinate chart in $U$, $K \subset V$ a compact set.
	We still denote it by $K$ for simple.
	On $\scrA_X^{p,q}(U)$, we have the following collection of seminorms:
	given a compact subset $K$ with a coordinate $(V,z)$ and an integer $m \geq 0$, define
	\[ p_{K,m}(\alpha) \coloneqq \max_{I; J; |\nu| \leq m; x\in K} \bigl|D^\nu\alpha_{I,J}(x)\bigr|, \quad \text{ where }\alpha|_V=\sum_{|I|=p,\,|J|=q}\alpha_{I,J}\,\upd z_I\wedge \upd\overline{z}_J. \]
	Here $D^\nu$ denotes differentiation in the real coordinates underlying $z$.

	Given two compact subset $K \subset K'$ with coordinates, easily see that we have $p_{K,m} \leq C_m p_{K',m}$ for some constant $C_m$ and for all $m$.
	Hence to define the topology by these seminorms, we only need to choose a covering of $X$ by compact subset $K_i$ with coordinates and consider the collection $\{p_{K_i,m}\}_{i,m}$ of seminorms.
	By the second countability, we can choose $\{K_i\}$ to be countable.
	Hence these seminorms make $\scrA_X^{p,q}(U)$ a Fr\'echet space.
	See \cite[Chapter~I, \S 2]{Demailly-2012-ComplexAnalyticDifferential}.

	Under this topology, we have $\alpha_n \to \alpha$ in $\scrA^{p,q}_X(U)$ iff for any compact subset $K$ with coordinates,
	the arbitrary derivatives of coefficient functions of $\alpha_n$ uniformly converge to the corresponding one of $\alpha$ on $K$.
\end{construction}


\begin{remark}[label=\labelNFM{rmk:recall-Frechet-space}, status=draft, tag=0000000]
	Here we recall some basic facts about Fr\'echet spaces.

	Recall that a \emph{seminorm} on a complex vector space $E$ is a map $p:E\to\bbR_{\geq0}$ such that
	\[
		p(\lambda v)=|\lambda|p(v),\qquad p(v+w)\leq p(v)+p(w).
	\]
	In other words, a semi-norm is a norm on $E/W$ for some subspace $W \subset E$.
	A collection of seminorms $(p_i)_{i \in I}$ can define a topology on $E$ by its weak topology.
	This topology is Hausdorff iff $(p_i)_{i \in I}$ separates points, i.e., for any $x,y \in E$, there exist $p_i$ with $p_i(x) \neq p_i(y)$.
	In particular, suppose the topology is Hausdorff, if $I$ is countable, then this topology can be induced by a metric
	\[
		\rho(v,w)=\sum_{j\geq1}2^{-j} \frac{p_j(v-w)}{1+p_j(v-w)}.
	\]
	And if $I$ if finite and the topology is Hausdorff, it can further be induced by a norm
	\[ \|v\| \coloneqq \sum_{1\leq j\leq N}p_j(v). \]

	A \emph{Fr\'echet space} is a complete Hausdorff locally convex topological vector space whose topology is defined by a countable family of seminorms $(p_j)_{j\geq1}$.
	A Fr\'echet space whose topology is defined by finitely many seminorms $p_1,\ldots,p_N$ is a Banach space.
\end{remark}

Let us record the basic properties of forms on a complex analytic variety.

\begin{proposition}[label=\labelNFM{prop:basic-properties-of-forms-with-topology}, status=draft, tag=0000000]
	Let $X$ be a complex analytic space of complex dimension $d$ and $U\subset X$ an open set.
	\begin{enumerate}
		\item The forms form a sheaf of rings.
		      More precisely, the wedge product makes $\scrA_X$ a sheaf of bi-graded topological $\cinf_{X,\bbC}$-algebras:
		      \[
			      \alpha\wedge\beta=(-1)^{(p+q)(r+s)}\beta\wedge\alpha,
			      \qquad \alpha\in\scrA_X^{p,q},\quad\beta\in\scrA_X^{r,s}.
		      \]
		      In particular, the subring $\bigoplus_p\scrA^{p,p}$ is commutative.
		\item There is a functorial continuous differential map
		      \[
			      \upd=\partial+\partialbar:\scrA_X^k\longrightarrow\scrA_X^{k+1},
			      \qquad
			      \partial : \scrA_X^{p,q}\longrightarrow\scrA_X^{p+1,q},\quad
			      \partialbar : \scrA_X^{p,q}\longrightarrow\scrA_X^{p,q+1}.
		      \]
		      They satisfy the Lebubinz rule and $d^2=\partial^2=\partialbar^2=0$ and
		      $\partial\partialbar+\partialbar\partial=0$.
		\item The integration gives a continuous evaluation map
		      \[
			      \int_U:\Gamma_c(\scrA_X^{d,d},U)\longrightarrow\bbC,
			      \qquad \alpha\longmapsto\int_U\alpha.
		      \]
		      In particular, if $X$ is compact, we have $\int: \scrA^{d,d}_X(X) \to \bbC$.
		\item Every holomorphic map $f:Y \to X$ defines a functorial continuous pullback $f^*: \scrA_X^{\bullet,\bullet}(U) \to \scrA_Y^{\bullet,\bullet}(f^{-1}(U))$ which is a homomorphism of bi-graded $\cinf_X$-algebras and commutative with $\partial,\partialbar$.
		\item There is a $\bbR$-linear map $\scrA_X \to \scrA_X, \alpha \mapsto \bar{\alpha}$.
		      The conjugation invariant forms form a subring of $\scrA_X$ (resp. $\bigoplus_p \scrA_X^{p,p}$), denoted by $\scrA_{X,\bbR}$ (resp. $\bigoplus_p \scrA_{X,\bbR}^{p,p}$).
		      The $\upd,\upd^c,f^*,\int_U$ commutate with the conjugation.
	\end{enumerate}
\end{proposition}

Philosophically, the forms serves as cocycles and now we introduce the \emph{currents}, which are topological dual of forms and serves as cycles.

\begin{construction}[label=\labelNFM{constr:currents}, status=draft, tag=0000000]
	Let $X$ be a complex analytic space of complex dimension $d$, $U\subset X$ an open set, and $0\leq p,q\leq d$.
	Set
	\[ \Gamma_c(\scrA_X^{p,q},U) =\varinjlim_{K\subset U \text{ compact}}\scrA_{X,K}^{p,q}(U),
		\qquad \scrA_{X,K}^{p,q}(U) \coloneqq \{\beta\in\scrA_X^{p,q}(U)\mid\Supp\beta\subset K\}. \]
	On each $K \subset U$ compact, we use the subspace topology $\scrA_{X,K}^{p,q}(U) \subset \scrA_{X}^{p,q}(U)$.
	And on $\Gamma_c(\scrA_X^{p,q},U)$, we use the inductive limit topology, i.e., a map $\Gamma_c(\scrA_X^{p,q},U) \to S$ is continuous iff $\scrA_{X,K}^{p,q}(U) \injmap \Gamma_c(\scrA_X^{p,q},U) \to S$ is continuous for all $K \subset U$ compact.

	Define the \emph{space of currents} $\scrA_{p,q}^X(U)$ to be the topological dual of $\Gamma_c(\scrA_X^{p,q},U)$, i.e.,
	\[ \scrA_{p,q}^X(U) \coloneqq \left\{\text{ linear and continuous map } T: \Gamma_c(\scrA_X^{p,q},U) \to \bbC \right\}; \]
	and an element in $\scrA_{p,q}^X(U)$ is called \emph{a currents of bidimension $(p,q)$} (or bidegree $(d-p,d-q)$.
	We equip $\scrA_{p,q}^X(U)$ with the weak topology defined by the seminorms $\{p_\alpha\}_{\alpha \in \Gamma_c(\scrA_X^{p,q},U)}$, where $p_\alpha(T) \coloneqq |T(\alpha)|$ with $T \in \scrA_{p,q}^X(U)$.

	For any $U \subset V$, the natural map is extension by zero:
	\[ i_*: \Gamma_c(\scrA_X^{p,q},U) \to \Gamma_c(\scrA_X^{p,q},V), \quad (i_*\alpha)|_{V\setminus U} = 0. \]
	This map is continuous.
	Then the dual of zero-extension $i_*$ gives a functorial restriction $\res^V_U: \scrA_{p,q}^X(V) \to \scrA_{p,q}^X(U)$.
	By partition of unity, the compatible local currents can be glued uniquely.
	Then we get a sheaf $\scrA_{p,q}^X$, called \emph{the sheaf of currents of bidimension $(p,q)$}.
	Similarly, we can define $\scrA_k^X$ and $\scrA_\bullet^X$ and we have
	\[ \scrA_k^X = \bigoplus_{p+q=k} \scrA_{p,q}^X, \qquad \scrA^X_\bullet = \bigoplus_k \scrA_k^X = \bigoplus_{p,q} \scrA_{p,q}^X. \]
	See \cite[Chapter~I, \S 2 and \S 3]{Demailly-2012-ComplexAnalyticDifferential}.
\end{construction}

\begin{remark}[label=\labelNFM{rmk:inductive-topo-on-the-space-for-compact-support-forms}, status=draft, tag=0000000]
	Note that in general, this topology on $\Gamma_c(\scrA_X^{p,q},U)$ is strictly finer than the subspace topology via $\Gamma_c(\scrA_X^{p,q},U) \subset \scrA_X^{p,q}(U)$.
	In other words, on $\Gamma_c(\scrA_X^{p,q},U)$, we have more open subsets, more continuous functions, and less convergance.

	For example, we consider in the $\cinf$ settings, set $U = \bbR$ be the real line as a real manifold.
	Let $f \in \cinf(\bbR)$ with $\Supp f \subset (0,1)$, $\int_\bbR f \upd x > 0$ and $f_n(x) \coloneqq f(x+n)$.
	With respect to the subspace topology, we have $f_n \to 0$.
	Consider the continuous function $T: \Gamma_c(\scrA_X^{0},U) \to \bbR, g \mapsto \int_\bbR g \upd x$ with respect to the inductive limit topology.
	We have $T(f_n) = T(f) \neq 0$ for all $n \geq 0$.
	Hence with respect to the inductive limit topology, we have NO the convergance $f_n \to 0$.
\end{remark}

Let $X$ be complex analytic space of purely complex dimension $d$.
Many object can be viewed as a currents.

\begin{example}[label=\labelNFM{eg:forms-as-currents}, status=draft, tag=0000000, title={Forms as currents}]
	Any $(p,q)$-forms $\alpha \in \scrA_X^{p,q}(U)$ can be viewed as a linear map by the integral
	\[ T_\alpha = \alpha \wedge - : \Gamma_c(\scrA_X^{d-p,d-q},U) \to \bbC, \quad \beta \mapsto \int_U \alpha \wedge \beta. \]
	Since the uniformly convergance commutates with integration, we get a $(d-p,d-q)$-currents.

	The inclusion $\scrA_X^{p,q} \to \scrA^X_{d-p,d-q}$ is continuous, but in general has no closed image.
	For example, let $f_n \in \cinf(\bbR)$ with $f_n(0) = 1, \Supp f_n \subset (-1/n,1/n)$.
	On $\scrA^1(\bbR)$, $\{f_n \upd x\}$ has no limit.
	But on $\scrA_0(\bbR)$, $f_n \upd x \to \delta_0$.
	This inclusion may be viewed as some analogue of $\scrC([0,1],\bbR) \injmap L^1([0,1],\bbR)$.

	By this reason, for a currents of bidimension $(p,q)$, we also call it has bidegree $(d-p,d-q)$.
\end{example}

\begin{example}[label=\labelNFM{eg:currents-of-dimension-2d}, status=draft, tag=0000000, title={Functions  as currents of top dimension}]
	Let us consider the currents of dimension $(d,d)$.
	By \cref{eg:forms-as-currents}, any $(0,0)$-forms, that is, $\cinf$-functions, can be viewed as a currents of bidimension $(d,d)$.
	As the construction in \cref{eg:forms-as-currents}, we only need to do integration.
	Hence, any locally integrable functions $f$ define a currents of bidimension $(d,d)$ or degree $0,0$.
\end{example}

\begin{example}[label=\labelNFM{eg:submainfolds-as-currents}, status=draft, tag=0000000, title={Cycles as currents}]
	Let $Z$ be a closed oriented real submanifold of $X$ of real dimension $k$ with piecewise regular boundary.
	For every open set $U\subset X$, set
	\[
		\delta_Z:\Gamma_c(\scrA_X^k,U)\longrightarrow\bbC,
		\qquad \beta\longmapsto\int_{Z\cap U}\beta|_Z.
	\]
	Since $Z$ is closed in $X$, the restriction of a compactly supported test form has compact support on $Z$.
	This define a currents $\delta_Z$ on $X$, called the \emph{delta current} associated to $Z$.
	By the Stokes theorem, $\upd \delta_Z = \delta_{\partial Z}$.

	If $Z$ is a closed complex subvariety of $X$ of complex dimension $p$ (or real dimension $k=2p$),
	we use the complex orientation and, if $Z$ is singular, define integration over its regular locus:
	\[
		\langle\delta_Z,\beta\rangle
		=\int_{Z_{\reg}\cap U}\beta|_{Z_{\reg}}.
	\]
	This defines a $k$-currents; for singular case, see \cite[Theorem 3.2.26]{Barlet-Magnusson-2019-ComplexAnalyticCycles}.
	Among test forms of total degree $k=2p$, only type $(p,p)$ can contribute: if $r+s=2p$ and $(r,s)\neq(p,p)$, then either $r>p$ or $s>p$, so the restriction to $Z_{\reg}$ vanishes.
	In other words, $\delta_Z \in \scrA_{p,p}^X$.
	In fact, $\delta_Z$ is a closed positive current, including when $Z$ is singular; see \cite[Chapter~III, Theorem~2.7]{Demailly-2012-ComplexAnalyticDifferential}.
\end{example}

\begin{example}[label=\labelNFM{eg:measure-and-distribution-as-currents}, status=draft, tag=0000000, title={Currents of dimension zero}]
	Let $\mu$ be any measure on the boral algebras with $\mu(K) < +\infty$ for any $K$ compact.
	Then
	\[ f \mapsto \int_X f \mu \]
	making $\mu$ a $0$-current.

	On $\scrA_{0}^X(X)$, there are more elements than measures since the topology on $\scrA_X^0(X)$ also capture the information of derivatives of the smooth functions.
	For example, fix a regular point $p \in X$ and a derivatives $D \in T_pX$ (i.e., a $\bbC$-linear map $(\cinf_{X,\bbC})_p \to \bbC$ satisfies the Lebubinz's rule), the linear $f \mapsto D(f)$ is a $0$-currents.
	It is not a measure since we can find a sequence $\{f_n\}$ of functions on a compact neighborhood of $p$ such that $f_n \to 0$ uniformly and $D(f_n) = 1$ for all $n$.
	For the $0$-currents, we have another name, \emph{distribution}.
	In the view point of above, the notion of distribution generalize the notion of measures.

	For simple, assume that
	\[ X = B(0,\underline{1}) = \left\{ (z_1,\cdots,z_d) \in \bbC^d \;\middle|\; |z_i| < 1 \right\}  \subset \bbC^d \]
	is the unit polydisc with coordinates $z_1,\cdots,z_d$.
	Given a $(p,q)$-currents $T$, for any $|I|=d-p,|J|=d-q$, we can define a distribution as
	\[ T_{I,J} : \Gamma_c(\cinf_{X,\bbC},X) \to \bbC, \quad f \mapsto \left< T, f \upd z^{I'}\wedge \upd \bar{z}^{J'}\right>, \]
	where $I \sqcup I' = J \sqcup J' = \{1,\cdots,d\} \coloneqq [d]$ and $z^{I'}\wedge \upd \bar{z}^{J'}$ has a ``suitable'' order.

	Here we imagine that $T_{I,J}$ is given by an ``integration'' $f \mapsto \int_X f\cdot T_{I,J} \upd z^{[d]} \wedge \upd \bar{z}^{[d]}$.
	Then formally consider the ``forms with distribution coefficients'', $\sum_{I,J} T_{I,J} \upd z^I \wedge \upd \bar{z}^J$, we have
	\begin{align*}
		\left< \sum_{I,J} T_{I,J} \upd z^I \wedge \upd \bar{z}^J, \sum_{I',J'} f_{I',J'}\upd z^{I'}\wedge \upd \bar{z}^{J'}\right> & = \int_X \left( \sum_{I,J,I',J'} T_{I,J} \cdot f_{I',J'} \upd z^{I\cup I'} \wedge \upd \bar{z}^{J\cup J'} \right) = \sum_{I,J} T_{I,J}(f_{[n]\setminus I,[n]\setminus J}) \\
		                                                                                                                           & = \sum_{I',J'} T(f_{I',J'}\upd z^{I'}\wedge \upd \bar{z}^{J'}) = \left<T, \sum_{I',J'} f_{I',J'}\upd z^{I'}\wedge \upd \bar{z}^{J'}\right>.
	\end{align*}
	The first line in the equation are some ``imagination'' without a real mathematical means.
	Hence formally we have $T = \sum_{I,J} T_{I,J} \upd z^I \wedge \upd \bar{z}^J$.
	From this view point, the distribution generalize the notion of functions and currents of dimension $(p,q)$ generalize the notion of forms of degree $(d-p,d-q)$.

	Historically, Schwarz first propose the notion of distribution, and then Georges de Rham introduce in notion of currents generalize distribution and forms; see \cite{Dinh-Sibony-2005-IntroductionTheoryCurrents}.
\end{example}

\begin{slogan}
	There is a philosophy analogue: forms are cocyles, like Cartier divisor or line bundles; and currents are cycles, like Weil divisor or subvarieties.
\end{slogan}

Let us record some basic properties of currents on a complex analytic spaces.

\begin{proposition}[label=\labelNFM{prop:basic-properties}, status=draft, tag=0000000]
	Let $X$ be a complex analytic variety of complex dimension $d$.
	We have the following properties about the currents.
	\begin{enumerate}
		\item The currents $\bigoplus \scrA^X_{p,q}$ forms a bi-graded topological $\scrA_X^{\bullet,\bullet}$-modules by the multiplication
		      \[ T \wedge \alpha \coloneqq \bigl( \beta \mapsto T(\alpha\wedge\beta) \bigr), \qquad T \in \scrA^X_{p,q}(U),\ \alpha \in \scrA_X^{r,s}(U),\ \beta \in \Gamma_c(\scrA_X^{p-r,q-s},U). \]
		      We also define $\alpha \wedge T \coloneqq (-1)^{(p+q)(r+s)}(T \wedge \alpha) \in \scrA^X_{p-r,q-s}(U)$ with $T \in \scrA^X_{p,q}(U),\alpha \in \scrA_X^{r,s}(U)$.
		      And it induces an inclusion
		      \[
			      \scrA_X^{p,q}\hookrightarrow\scrA_{d-p,d-q}^X, \qquad \alpha\longmapsto T_\alpha, \qquad \langle T_\alpha,\beta\rangle=\int_U \alpha \wedge \beta.
		      \]
		\item The differential on $\scrA^{\bullet,\bullet}$ gives differentials on $\scrA_{\bullet,\bullet}$ by
		      \[
			      \langle\partial T,u\rangle\coloneqq(-1)^{p+q+1}\langle T,\partial u\rangle,
			      \qquad
			      \langle\partialbar T,u\rangle\coloneqq(-1)^{p+q+1}\langle T,\partialbar u\rangle,
		      \]
		      where $T\in\scrA_{p,q}^X(U)$ and $u$ is a test form of the appropriate type.
		      And we set $\upd \coloneqq \partial + \partialbar$ and $\upd^c\coloneqq (\partial-\partialbar)/(2\pi \upi)$.
		      Then we have
		      \[
			      \upd, \upd^c:\scrA_k^X \longrightarrow\scrA^X_{k-1},
			      \qquad
			      \partial : \scrA^X_{p,q}\longrightarrow\scrA^X_{p-1,q},\quad
			      \partialbar : \scrA^X_{p,q}\longrightarrow\scrA^X_{p,q-1}.
		      \]
		      and
		      $\upd^2=\partial^2=\partialbar^2=0$,
		      $\partial\partialbar+\partialbar\partial=0$, and
		      $\upd T_\alpha=T_{\upd\alpha}$.
		\item For a holomorphic submersion $\pi:E \to X$ of complex manifolds of relative complex dimension $r$, there is a functorial pullback
		      \[
			      \pi^*:\scrA_{p,q}^X(U)\longrightarrow\scrA_{p+r,q+r}^E(\pi^{-1}(U)),
			      \qquad \langle \pi^*T,\beta\rangle=\langle T,\pi_!\beta\rangle,
		      \]
		      where $\pi_!\beta$ is the smooth compactly supported $(p,q)$-form obtained by integrating a test $(p+r,q+r)$-form $\beta$ along the fibers, with their complex orientation.
					We have 
					\[ \partial\pi^*T=\pi^*\partial T,\qquad \partialbar\pi^*T=\pi^*\partialbar T \qquad \pi^*T_{\alpha} = T_{\pi^*\alpha}. \]
					See \cite[Proposition 3.2.5]{Dinh-Sibony-2005-IntroductionTheoryCurrents}.
		      % \par\latexmath{Here $E$ and $X$ are complex manifolds.}
		      % \latexdeleted{\Yang{commutative}}
		      % \[
			    %   \latexaddedmath{\partial\pi^*T=\pi^*\partial T,\qquad
				  %     \partialbar\pi^*T=\pi^*\partialbar T.}
		      % \]
		      % And we have $\pi^*T_{\alpha} = T_{\pi^*\alpha}$.
		      % \latexdeleted{\Yang{ref}}
		      % \begin{latexmathchange}
			    %   See \cite[Chapter~I, \S 2.C.1--2.C.2, (2.15)--(2.18)]{Demailly-2012-ComplexAnalyticDifferential}.
			    %   No properness is needed here: $\Supp(\pi_!\beta)\subset\pi(\Supp\beta)$ is compact for every compactly supported test form $\beta$.
		      % \end{latexmathchange}
		      % \toai{find a references} -- resolved in place.
		\item For a proper holomorphic map $f: X \to Y$ between complex manifolds, there is a functorial pushforward
		      \[
			      f_*:\scrA_{p,q}^X(X)\longrightarrow\scrA_{p,q}^Y(Y),\qquad \langle f_*T,\beta\rangle \coloneqq \langle T,f^*\beta\rangle.
		      \]
		      The pushforward is a continuous homomorphism of bi-graded $\scrA_Y^{\bullet,\bullet}$-modules.
		      In particular, we have the projection formula
		      \begin{align*}
			      f_*(T \wedge f^*\alpha)&=(f_* T) \wedge \alpha
			      \in\latexdeletedmath{...}\latexaddedmath{\scrA_{p-r,q-s}^Y(Y)},\\
			      \alpha&\in\latexaddedmath{\scrA_Y^{r,s}(Y)},\qquad
			      T\in\latexaddedmath{\scrA_{p,q}^X(X)}.
		      \end{align*}
		      \latexmath{Here $0\leq r\leq p$ and $0\leq s\leq q$.}
					% \toai{complete the formula} -- resolved in place.
					See \cite[Exercise 3.2.12]{Dinh-Sibony-2005-IntroductionTheoryCurrents}.
		      % \latexdeleted{\Yang{ref}}
		      % \begin{latexmathchange}
			    %   The source module structure is induced by $f^*$.
			    %   See \cite[Chapter~I, \S 2.C.1, Theorem~2.14]{Demailly-2012-ComplexAnalyticDifferential}.
			    %   More generally, $f_*T$ is defined when $f|_{\Supp T}$ is proper.
			    %   In this case, pair $T$ with $\chi f^*\beta$, where $\chi\in\cinf_c(X)$ equals $1$ near the compact set $\Supp T\cap f^{-1}(\Supp\beta)$; the result is independent of $\chi$.
			    %   The support condition is needed because $f^*\beta$ need not have compact support.
		      % \end{latexmathchange}
		      % % \toai{find a references} -- resolved in place.
		\item The conjugation on $\scrA^{\bullet}$ induces conjugation on $\scrA_{\bullet}$ by $\left< \bar{T}, \alpha \right> \coloneqq \overline{\left< T, \bar{\alpha} \right>}$.
		      As forms, the conjugation invariant currents form an $\scrA_{X,\bbR}^{\bullet}$-module $\scrA^{X,\bbR}_{\bullet}$ (resp. $\bigoplus_p \scrA_{X,\bbR}^{p,p}$-module $\bigoplus_p \scrA^{X,\bbR}_{p,p}$).
		      The $\upd,\upd^c,\pi^*,f_*$ commutate with the conjugation.
	\end{enumerate}
	\Yang{}
\end{proposition}

We now compare the homology of currents with the cohomology of forms.
First recall the definitions.

\begin{construction}[label=\labelNFM{constr:cohomology-by-forms-and-currents}, status=draft, tag=0000000]
	Let $X$ be a complex manifold of complex dimension $d$.
	We have the following complexes:
	\begin{align*}
		\calC_{\dR, \bullet}
		 & : 0 \longrightarrow \scrA_X^0 \xrightarrow{\upd}\scrA_X^1 \xrightarrow{\upd}\cdots
		\xrightarrow{\upd}\scrA_X^{2d}
		\latexdeletedmath{\xrightarrow{\upd}\cdots}\latexaddedmath{\longrightarrow0},                            \\
		\calC_{\dR, \bbR, \bullet}
		 & : \latexdeletedmath{0 \longrightarrow \scrA_X^0 \xrightarrow{\upd}\scrA_X^1
		\xrightarrow{\upd}\cdots\xrightarrow{\upd}\scrA_X^{2d}\xrightarrow{\upd}\cdots},                         \\[-2pt]
		 & \phantom{:}\latexaddedmath{0\longrightarrow\scrA_{X,\bbR}^0\xrightarrow{\upd}\scrA_{X,\bbR}^1
		\xrightarrow{\upd}\cdots\xrightarrow{\upd}\scrA_{X,\bbR}^{2d}\longrightarrow0},                          \\
		\calC_{\Db,\bullet}^{p}
		 & : 0\longrightarrow\scrA_X^{p,0}\xrightarrow{\partialbar}\scrA_X^{p,1}
		\xrightarrow{\partialbar}\cdots\xrightarrow{\partialbar}\scrA_X^{p,d}
		\latexdeletedmath{\xrightarrow{\partialbar}\cdots}\latexaddedmath{\longrightarrow0},                     \\
		\calC^{\mathrm{cur}}_{\bullet}
		 & : \latexdeletedmath{\cdots}\latexaddedmath{0}\xrightarrow{\upd}\scrA_{2d}^X
		\xrightarrow{\upd}\scrA_{2d-1}^X\xrightarrow{\upd}\cdots\xrightarrow{\upd}\scrA_0^X\longrightarrow0,     \\
		\calC^{\mathrm{cur},\bbR}_{\bullet}
		 & : \latexdeletedmath{\cdots\xrightarrow{\upd}\scrA_{2d}^X\xrightarrow{\upd}\scrA_{2d-1}^X
		\xrightarrow{\upd}\cdots\xrightarrow{\upd}\scrA_0^X\longrightarrow0},                                    \\[-2pt]
		 & \phantom{:}\latexaddedmath{0\longrightarrow\scrA_{2d}^{X,\bbR}\xrightarrow{\upd}\scrA_{2d-1}^{X,\bbR}
		\xrightarrow{\upd}\cdots\xrightarrow{\upd}\scrA_0^{X,\bbR}\longrightarrow0},                             \\
		\calC_{p,\bullet}^{\mathrm{cur}}
		 & : \latexdeletedmath{\cdots}\latexaddedmath{0}\xrightarrow{\partialbar}\scrA_{p,d}^X
		\xrightarrow{\partialbar}\scrA_{p,d-1}^X\xrightarrow{\partialbar}\cdots
		\xrightarrow{\partialbar}\scrA_{p,0}^X\longrightarrow0.
	\end{align*}
	\latexmath{The real terms are the conjugation invariant subsheaves.}
	% \toai{revise the complex with real coefficient} -- resolved in place.
	On each complex, \latexdeleted{the $0$ is put in the degree $0$ part.}
	\latexmath{the term with index $0$ is in homological degree $0$.}
	\begin{latexmathchange}
		More precisely, $\scrA_X^k$ is in degree $-k$ and $\scrA_k^X$ in degree $k$.
		In the Dolbeault complexes, $\scrA_X^{p,q}$ is in degree $-q$ and $\scrA_{p,q}^X$ in degree $q$.
		The same conventions apply to real coefficients.
	\end{latexmathchange}
	Recall that for a complex $\calC_\bullet$, its homological groups are defined by
	\[
		H_n(\calC_\bullet)\coloneqq
		\latexdeletedmath{\Ker\upd_k/\Image\upd_{k+1}}
		\latexaddedmath{\Ker\upd_n/\Image\upd_{n+1}}.
	\]
	Taking the homological groups of the \latexdeleted{four}\latexmath{six} \latexdeleted{complex}\latexgrammar{complexes}, we have
	\begin{align*}
		H^k_\dR(X,\bbC)
		 & \coloneqq H_{-k}\bigl(\Gamma(X,\calC_{\dR,\bullet})\bigr),                                        \\
		H_{\dR}^k(X,\bbR)
		 & \latexaddedmath{\coloneqq H_{-k}\bigl(\Gamma(X,\calC_{\dR,\bbR,\bullet})\bigr)},                  \\
		H^{p,q}_{\partialbar}(X,\bbC)
		 & \coloneqq H_{-q}\bigl(\Gamma(X,\calC_{\Db,\bullet}^{p})\bigr),                                    \\
		H_k^{\mathrm{cur}}(X,\bbC)
		 & \coloneqq H_k\bigl(\Gamma(X,
		\latexdeletedmath{\calC_{\mathrm{cur},\bullet}}\latexaddedmath{\calC^{\mathrm{cur}}_\bullet})\bigr), \\
		H^{cur}_k(X,\bbR)
		 & \latexaddedmath{\coloneqq H_k\bigl(\Gamma(X,\calC^{\mathrm{cur},\bbR}_\bullet)\bigr)},            \\
		H_{p,q}^{\partialbar,\mathrm{cur}}(X,\bbC)
		 & \coloneqq H_q\bigl(\Gamma(X,
			\latexdeletedmath{\calC_{\mathrm{cur},\bullet}^{p}}\latexaddedmath{\calC_{p,\bullet}^{\mathrm{cur}}})\bigr).
	\end{align*}
	% \toai{finish the definition of real coefficient group} -- resolved in place.
	As the case of forms, we say that a currents $T$ is \emph{closed} (resp, \emph{exact}, \emph{$\partialbar$-closed}, \emph{$\partialbar$-exact}) if $T \in \Ker \upd$ (resp, $T \in \Image \upd$, $T \in \Ker \partialbar$, $T \in \Image \partialbar$).
\end{construction}


\begin{theorem}[label=\labelNFM{thm:cohomology-by-forms-and-currents}, status=draft, tag=0000000]
	Let $X$ be a complex manifold of complex dimension $d$.
	\begin{enumerate}
		\item (Comparison) The inclusions $\scrA^{p,q} \to \scrA_{d-p,d-q}$ give isomorphisms
		      \[ H^{k}_\dR(X,\bbC) \cong H_{\latexdeletedmath{d-k}\latexaddedmath{2d-k}}^{cur}(X,\bbC), \qquad H^{p,q}_{\partialbar}(X, \bbC) \cong H^{\partialbar,cur}_{d-p,d-q}(X,\bbC); \]
		      the acyclic resoluations give isomorphisms
		      \[ H^{k}_\dR(X,\bbR) \cong H^k(X,\underline{\bbR}), \quad H^{\latexdeletedmath{p}\latexaddedmath{q}}(X,\Omega_X^{\latexdeletedmath{q}\latexaddedmath{p}}) \cong H^{p,q}_{\partialbar}(X, \underline{\bbC}). \]
		      When $X$ is compact, the map $Z \mapsto \delta_Z$ gives isomorphisms 
		      \[ H_k^{cur}(X,\bbR) \xleftarrow{\cong} H_k^{\sing,manifold}(X,\bbR) \xrightarrow{\cong} H_{k}^{\sing}(X,\bbR); \]
		\item ($\ppbar$-lemma) Suppose that $X$ is a compact K\"ahler manifold, and $\alpha$ is a $\upd$-closed $(p,q)$-form.
		      Then $\alpha$ is $\upd$-exact iff it is $\partial$-exact iff it is $\partialbar$-exact iff it is $\ddc$-exact.
		\item (Hodge decomposition) Suppose that $X$ is a compact K\"ahler manifold.
		      Then the decompositions of $\cinf$-forms give decompositions
		      \[ H^k_{\dR}(X,\bbC) = \bigoplus_{p+q=k} H^{p,q}_{\partialbar}(X,\bbC), \quad H^{p,q}_{\partialbar}(X,\bbC) \cong \overline{H^{q,p}_{\partialbar}(X,\bbC)}. \]
	\end{enumerate}
	\Yang{}
\end{theorem}
\begin{proof}[References and sketch of proof]

	Here we give a sketch of (c).

	for the ddc lemma, see \cite[Proposition 6.17]{Voisin-2002-HodgeTheoryComplex}.
	\Yang{Demailly IV-6}
	\Yang{}

\end{proof}


\begin{latexmathchange}
	\textbf{Proposed mathematical revision of the comparison statement (a).}
	Let $X$ be a complex manifold of complex dimension $d$.
	The inclusions of forms into currents induce
	\[
		H^k_\dR(X,\bbC)\cong H_{2d-k}^{\mathrm{cur}}(X,\bbC),\qquad
		H^{p,q}_{\partialbar}(X,\bbC)\cong H_{d-p,d-q}^{\partialbar,\mathrm{cur}}(X,\bbC).
	\]
	The first index is $2d-k$ because $X$ has real dimension $2d$.
	The de Rham and Dolbeault resolutions give
	\[
		H^k_\dR(X,\bbR)\cong H^k(X,\underline{\bbR}),\qquad
		H^{p,q}_{\partialbar}(X,\bbC)\cong H^q(X,\Omega_X^p).
	\]
	See \cite[Chapter~IV, \S 6]{Demailly-2012-ComplexAnalyticDifferential}.

	For the singular homology comparison, the support convention matters.
	Define compactly supported current homology by
	\[
		H_k^{\mathrm{cur},c}(X,\bbR)
		\coloneqq H_k\bigl(\Gamma_c(X,\calC^{\mathrm{cur},\bbR}_\bullet)\bigr).
	\]
	Integration over smooth singular simplices gives isomorphisms
	\[
		H_k^{\mathrm{cur},c}(X,\bbR)
		\xleftarrow{\cong} H_k^{\sing,\mathrm{sm}}(X,\bbR)
		\xrightarrow{\cong} H_k^{\sing}(X,\bbR),
	\]
	where the middle group is computed by finite chains of smooth singular simplices, not just embedded submanifolds.
	For the unrestricted currents used above, locally finite smooth singular chains instead give
	\[
		H_k^{\mathrm{cur}}(X,\bbR)\cong H_k^{\mathrm{BM}}(X,\bbR),
	\]
	the Borel--Moore homology of $X$.
	If $X$ is compact, the support distinctions disappear and the original singular homology comparison applies.
	See \cite[\S 2, especially the discussion after (2.8), and \S 3, Theorem~3.2]{Iversen-1989-CauchyResiduesDeRham}.

	There is also a sign to record: with the differential on currents chosen above, Stokes gives
	\[
		\upd\delta_\sigma=(-1)^{k+1}\delta_{\partial\sigma}
	\]
	for an oriented smooth $k$-simplex $\sigma$ (and similarly for a submanifold with boundary).
	Consequently the chain map for the comparison is
	$\sigma\mapsto(-1)^{k(k-1)/2}\delta_\sigma$.
	This sign adjustment preserves the chosen identity $\upd T_\alpha=T_{\upd\alpha}$.
\end{latexmathchange}

Now let us turn to the positivity.
Forms and currents are more finer object than algebraic cycles and cocyles.
They can capture the ``local positivity''.

First we record the positivity of real $(1,1)$-forms.

\begin{definition}[label=\labelNFM{def:positivity-of-11-forms}, status=draft, tag=0000000]
	Let $X$ be a complex manifold, $x \in X$ and $\alpha \in \scrA_{X,\bbR}^{1,1}(X)$.
	\par\latexmath{Let $J$ be the complex structure on $T_xX$.}
	We say that $\alpha$ is \emph{positive at $x$} if the \latexdeleted{bilinear}\latexmath{sesquilinear} map
	\[
		T_xX \times T_xX \to \bbC,\quad
		(u,v)\mapsto\alpha(u,Jv)
		\latexdeletedmath{+\upi}\latexaddedmath{-\upi\alpha(u,v)}
	\]
	is a hermitian inner product on \latexdeleted{$T_x$}\latexmath{$T_xX$}.
	\latexmath{This Hermitian form is linear in its first argument.}
	We  say that $\alpha$ is \emph{semi-positive at $x$} \latexdeleted{is limit of ...}
	\latexmath{if $\alpha_x(u,Ju)\geq0$ for every $u\in T_xX$.}
	\begin{latexmathchange}
		Equivalently, $\alpha_x$ is a limit of positive definite real $(1,1)$-forms on $T_xX$.
		Indeed, for any positive definite $\omega_x$, the forms $\alpha_x+\varepsilon\omega_x$ are positive definite for every $\varepsilon>0$ and converge to $\alpha_x$.
		The associated symmetric, $J$-invariant real bilinear form is $g_\alpha(u,v)=\alpha(u,Jv)$.
	\end{latexmathchange}
	We say \latexdeleted{the}\latexgrammar{that} $\alpha$ is \emph{positive} (resp semi-positive) on $X$ \latexdeleted{if on every point.}
	\par
	\latexgrammar{if it has this property at every point.}
	\begin{latexmathchange}
		Write $\alpha>0$ and $\alpha\geq0$, respectively.
		See \cite[Chapter~III, \S 1.A]{Demailly-2012-ComplexAnalyticDifferential}.
	\end{latexmathchange}
	% \toai{finish and polish the definition} -- resolved in place.
	\Yang{}
\end{definition}


\begin{lemma}[label=\labelNFM{lem:characterize-of-positive-11-form}, status=draft, tag=0000000]
	Let $X$ be a complex manifold, $x \in X$ and $\alpha \in \scrA_{X,\bbR}^{1,1}(X)$.
	\latexmath{Set $d=\dim_\bbC X\geq1$.}
	TFAE:
	\begin{enumerate}
		\item $\alpha$ is positive \latexmath{(resp. semi-positive)} at $x$;
		\item $\latexdeletedmath{\alpha = \upi \sum u_i \beta_i \wedge \bar{\beta}}$
		      \latexdeleted{for $\beta_i,\cdots,\beta_d \in \scrA^{1,0}_X$ a basis of $T_xX$.}
		      \[
			      \latexaddedmath{\alpha_x=\upi\sum_{j=1}^d\lambda_j\beta_j\wedge\bar\beta_j,
				      \qquad\lambda_j>0\quad\text{(resp. }\lambda_j\geq0\text{)},}
		      \]
		      \latexmath{where $\beta_1,\ldots,\beta_d$ is a coframe of $T_x^{*1,0}X$.}
		\item \latexdeleted{$\alpha \wedge \beta \geq 0$ for all beta.}
		      \begin{latexmathchange}
			      $\alpha_x\wedge\beta>0$ (resp. $\geq0$) for every nonzero strongly positive $(d-1,d-1)$-covector $\beta$ at $x$.
			      The inequality means a positive (resp. nonnegative) multiple of the complex orientation volume form.
			      It is enough to test $\beta=(\upi\eta_1\wedge\bar\eta_1)\wedge\cdots\wedge(\upi\eta_{d-1}\wedge\bar\eta_{d-1})$ with the $\eta_j$ linearly independent; for $d=1$, use $\beta=1$.
		      \end{latexmathchange}
		\item \latexmath{In holomorphic coordinates,}
		      \begin{latexmathchange}
			      $\alpha_x=\upi\sum_{j,k=1}^d a_{j\bar k}\,\upd z_j\wedge\upd\bar z_k$, where the Hermitian matrix $(a_{j\bar k})$ is positive definite (resp. positive semidefinite).
		      \end{latexmathchange}
	\end{enumerate}
	\begin{latexmathchange}
		Diagonalization gives the coframe criterion, and the complementary wedge tests detect each eigenvalue.
		See \cite[Chapter~III, \S 1.A, (1.8)--(1.9)]{Demailly-2012-ComplexAnalyticDifferential}.
	\end{latexmathchange}
	% \toai{finish and polish the definition} -- resolved in place.
	\Yang{}
\end{lemma}


\begin{definition}[label=\labelNFM{def:kalher-and-hermitian-manifold}, status=draft, tag=0000000]
	By a K\"ahler manifold (resp. \emph{Hermitian manifold}), we mean a pair $(X,\omega)$, where $X$ is a complex manifold and $\omega$ a real positive closed $(1,1)$-$\cinf$-forms (resp. a real positive $(1,1)$-$\cinf$-forms ).
	% \begin{latexmathchange}
	% 	Thus $\upd\omega=0$ is the additional K\"ahler condition; a Hermitian form need not be closed.
	% 	The associated Riemannian metric is $g_\omega(u,v)=\omega(u,Jv)$.
	% 	See \cite[Chapter~V, \S 7, and Chapter~VI, \S 5]{Demailly-2012-ComplexAnalyticDifferential}.
	% \end{latexmathchange}
	% % \toai{finish and polish the definition} -- resolved in place.
	\Yang{}
\end{definition}


\begin{definition}[label=\labelNFM{def:positivity-of-forms-and-currents}, status=draft, tag=0000000]
	Let $X$ be a complex manifold \latexmath{of complex dimension $d$} and $\gamma \in \scrA_{X,\bbR}^{p,p}(X)$ be a form , $T \in \scrA^{X,\bbR}_{p,p}$ be a \latexdeleted{currents.....}\latexgrammar{current.}
	\latexmath{Assume $0\leq p\leq d$; the form conditions are pointwise.}
	\begin{enumerate}
		\item \latexdeleted{We say the $\gamma$ is \emph{strongly positive} (resp. semi-strongly positive) if $\gamma = \sum a_i \alpha_{i_1} \wedge\cdots \wedge \alpha_{i,p}$ with ...}
		      \begin{latexmathchange}
			      We say that $\gamma$ is \emph{strongly positive} if, at every $x$,
			      \[
				      \gamma_x=\sum_a c_a
				      (\upi\eta_{a,1}\wedge\bar\eta_{a,1})\wedge\cdots
				      \wedge(\upi\eta_{a,p}\wedge\bar\eta_{a,p}),
				      \qquad c_a\geq0,
			      \]
			      with $\eta_{a,j}\in T_x^{*1,0}X$.
			      Equivalently, the generators are $\upi^{p^2}\eta\wedge\bar\eta$ with $\eta$ a decomposable $(p,0)$-covector.
		      \end{latexmathchange}
		\item \latexdeleted{We say that $\gamma$ is \emph{(weakly) positive} if $\gamma \cdot \beta > 0$ ($\geq 0$) for all semi-strongly positive $\beta$.}
		      \begin{latexmathchange}
			      We say that $\gamma$ is \emph{weakly positive}, or simply \emph{positive}, if $\gamma_x\wedge\beta\geq0$ for every strongly positive $(d-p,d-p)$-covector $\beta$ at $x$.
			      Here a top degree form is nonnegative if it is a nonnegative real multiple of the complex orientation volume form.
		      \end{latexmathchange}
		\item \latexmath{There is also an intermediate notion:}
		      \begin{latexmathchange}
			      $\gamma$ is \emph{Hermitian semi-positive} if
			      \[
				      \gamma_x=\upi^{p^2}\sum_a c_a\theta_a\wedge\bar\theta_a,
				      \qquad c_a\geq0,\quad\theta_a\in\bigwedge^pT_x^{*1,0}X,
			      \]
			      without requiring the $\theta_a$ to be decomposable.
			      Equivalently, the Hermitian coefficient matrix in
			      $\gamma_x=\upi^{p^2}\sum_{I,J}a_{I\bar J}\,\upd z_I\wedge\upd\bar z_J$ is positive semidefinite.
			      See \cite[\S 1, Definitions~1.1 and~1.4]{Debarre-Ein-Lazarsfeld-Voisin-2010-PseudoeffectiveNefClasses}.
		      \end{latexmathchange}
		\item \latexdeleted{We say that $T$ is \emph{strongly positive} (resp. \emph{(weakly) positive}) if $\left<T,\beta\right> > 0$ for all $\beta \in $.}
		      \begin{latexmathchange}
			      A real current $T$ of bidimension $(p,p)$ is \emph{positive} (or \emph{weakly positive}) if $\langle T,\beta\rangle\geq0$ for all smooth compactly supported strongly positive $(p,p)$-forms $\beta$.
			      It is \emph{strongly positive} if the same inequality holds for all smooth compactly supported weakly positive $(p,p)$-forms.
			      Testing against Hermitian semi-positive $(p,p)$-forms similarly defines Hermitian semi-positive currents.
			      The stronger condition uses the larger test cone.
			      The inequality is non-strict, since the zero test form is allowed and a nonzero test form may be supported away from $\Supp T$.
			      These definitions agree with the pointwise ones for smooth forms viewed as currents.
			      Integration over a closed complex $p$-dimensional subvariety gives a strongly positive current.
		      \end{latexmathchange}
	\end{enumerate}
	\begin{latexmathchange}
		For both forms and currents, strong positivity implies Hermitian semi-positivity, which implies weak positivity.
		For forms, the three closed cones coincide in degrees $p=0,1,d-1,d$.
		We use \emph{strictly} for their respective interiors; thus the draft's \emph{semi-strongly positive} cone is usually called strongly positive.
		In degree $(1,1)$, this nonnegative convention agrees with semi-positivity above, and its strict version agrees with positive definiteness.
		See \cite[Chapter~III, \S 1.A--1.B and \S 2]{Demailly-2012-ComplexAnalyticDifferential}.
	\end{latexmathchange}
	% \toai{is there other positivity?, add with reference, and polish the original definition} -- resolved in place.
	\Yang{}
\end{definition}

\begin{remark}[label=\labelNFM{rmk:comparing-with-algebraic-geometry}, status=draft, tag=0000000]

	The positivity in algebraic geometry

	\Yang{}
\end{remark}









\subsection[status=draft,tag=0000000,label=\labelNFM{subsec:potential-theory}]{Potential theory}

In this subsection, we work on a complex manifold $X$ of complex dimension $d$,
except when discussing psh functions on singular spaces.
We first describe potentials of currents, then introduce psh functions and use them to define products of currents.

Recall our convention in \cref{eg:measure-and-distribution-as-currents}:
a \emph{distribution} is a current of dimension zero, hence of bidegree $(d,d)$.
Its exterior differential is zero.
A $\ddc$-potential of a current $T$ is a current $U$ satisfying $T = \ddc U$.
For $T$ of bidegree $(1,1)$, the potential $U$ has bidegree $(0,0)$,
or equivalently bidimension $(d,d)$.
The operator on these potentials is
\[
    \ddc: \scrA_{d,d}^X \longrightarrow \scrA_{d - 1, d - 1}^X.
\]
As in \cref{eg:currents-of-dimension-2d}, a locally integrable function $u$ is viewed as a current of bidegree $(0,0)$ by
\[
    \langle u, \alpha \rangle \coloneqq \int_X u \alpha, \qquad \alpha \in \Gamma_c(\scrA_X^{d,d}, X).
\]

\begin{proposition}[label=\labelNFM{prop:local-ddc-potentials}, status=draft, tag=0000000]
    Let $T$ be a closed real current of bidegree $(1,1)$ on $X$.
    For every $x \in X$, there exist an open neighborhood $U$ and a real current $u$ of bidegree $(0,0)$ on $U$ such that
    \[
        T|_U = \ddc u = \frac{\im}{\pi} \partial \partialbar u.
    \]
    We call $u$ a \emph{local potential} of $T$.
    If $T$ is smooth, $u$ can be chosen smooth.

    Two real local potentials for the same current differ by a smooth \emph{pluriharmonic} function,
    i.e., a real function $h$ with $\partial \partialbar h = 0$.
    Such a function is locally the real part of a holomorphic function.
    See \cite[Chapter~III, \S 1.18 and Proposition~1.19]{Demailly-2012-ComplexAnalyticDifferential}.
\end{proposition}
\begin{proof}
    The assertion is local.
    By the Poincar\'e lemma for currents, after shrinking a neighborhood of $x$,
    we can write $T = \upd S$ with $S$ a real current of degree $1$.
    Write $S = S^{1,0} + S^{0,1}$.
    Since $\partialbar S^{0,1} = T^{0,2} = 0$,
    the Dolbeault--Grothendieck lemma gives $S^{0,1} = \partialbar v$ locally for a complex current $v$ of bidegree $(0,0)$.
    As $S$ is real, we have $S^{1,0} = \partial \overline{v}$, and hence
    \[
        T = \partial \partialbar (v - \overline{v}) = \ddc \bigl(2\pi \operatorname{Im} v\bigr).
    \]
    If $T$ is smooth, the smooth versions of these lemmas give a smooth potential.

    Finally, the difference $h$ of two potentials satisfies $\ddc h = 0$.
    Taking the trace in local coordinates gives $\Delta h = 0$ weakly.
    By elliptic regularity, $h$ is smooth, and then pluriharmonic.
\end{proof}

\begin{proposition}[label=\labelNFM{prop:global-ddc-potentials}, status=draft, tag=0000000]
    Suppose that $X$ is connected, compact and K\"ahler.
    Let $\omega$ be a smooth closed real $(1,1)$-form and $T$ a closed real current of bidegree $(1,1)$.
    Then there exists a real current $u$ of bidegree $(0,0)$ on $X$ such that
    \[
        T = \omega + \ddc u
    \]
    iff $[T] = [\omega]$ in $H^2_{\dR}(X, \bbR)$.
    Such a potential is unique up to an additive real constant.
\end{proposition}
\begin{proof}
    If $T = \omega + \ddc u$, then $T - \omega = \upd(\upd^c u)$ is exact,
    so $[T] = [\omega]$.
    Conversely, apply the $\partial \partialbar$-lemma for currents to $T - \omega$.
    The Hodge decomposition proof of the smooth version in \cref{thm:cohomology-by-forms-and-currents} also applies to currents,
    and gives a real current $u$ of bidegree $(0,0)$ with $T - \omega = \ddc u$.
    See \cite[Chapter~VI, \S 3 and Lemma~8.6]{Demailly-2012-ComplexAnalyticDifferential} for the argument using Hodge decomposition.

    The difference of two global potentials is a smooth pluriharmonic function on a connected compact manifold,
    hence constant by the maximum principle.
\end{proof}

We now characterize the local potentials of positive currents.

\begin{definition}[label=\labelNFM{def:psh-function}, status=draft, tag=0000000]
    Let $u: X \to [-\infty, +\infty)$ be a function.
    We say that $u$ is \emph{plurisubharmonic}, or \emph{psh} for short,
    if it satisfies the following conditions:
    \begin{enumerate}
        \item $u$ is upper semicontinuous, i.e., $u^{-1}([-\infty, a)) = \{x \in X \mid u(x) < a\}$ is open for every $a \in \bbR$.
        \item For every holomorphic map $\gamma: \Delta \to X$ from the unit disc $\Delta \subset \bbC$,
              the function $u \circ \gamma$ is subharmonic or identically $-\infty$.
    \end{enumerate}
    Denote the set of psh functions on $X$ by $\operatorname{PSH}(X)$.
    Here we allow the function identically $-\infty$.
    For an upper semicontinuous function, the second condition can also be checked on complex affine lines in each coordinate chart.
    See \cite[Chapter~I, \S 5]{Demailly-2012-ComplexAnalyticDifferential}.
\end{definition}

\begin{remark}[label=\labelNFM{rmk:psh-function-in-real-and-complex-one-dimensional-and-differentiable-case}, status=draft, tag=0000000]
    We compare this definition with some familiar cases.
    \begin{enumerate}
        \item The real one-dimensional analogue is convexity.
              More precisely, a real-valued function $v$ on an open interval $I \subset \bbR$ is convex iff $u(z) = v(\operatorname{Re} z)$ is psh on $I + \im \bbR$.
        \item On an open subset of $\bbC$, psh is the same as subharmonic,
              with the same convention for the function identically $-\infty$.
              For a $\scrC^2$ real-valued function, this is equivalent to $\Delta u \geq 0$.
        \item If $u \in \scrC^2(X, \bbR)$, then $u$ is psh iff its Levi matrix
              \[
                  \left(\frac{\partial^2 u}{\partial z_j \partial \overline{z}_k}\right)_{1 \leq j,k \leq d}
              \]
              is positive semidefinite in every coordinate chart.
              Equivalently, the real $(1,1)$-form $\ddc u = \frac{\im}{\pi} \partial \partialbar u$ is positive.
    \end{enumerate}
    More generally, if $u$ is psh and is not identically $-\infty$ on any connected component,
    then $u \in L^1_{\mathrm{loc}}(X)$ and $\ddc u$ is a closed positive current of bidegree $(1,1)$.
    Conversely, a real current $U$ of bidegree $(0,0)$ with $\ddc U \geq 0$ is represented by a unique psh function.
    Thus a closed positive current of bidegree $(1,1)$ has local potentials given by psh functions.
    See \cite[Chapter~I, \S 5, especially Theorem~5.8, and Chapter~III, Proposition~1.19]{Demailly-2012-ComplexAnalyticDifferential}.
\end{remark}

On a reduced complex analytic space, we say that $u$ is psh if,
locally under a closed embedding into a complex manifold,
it is the restriction of a psh function on an ambient neighborhood.
This extends the preceding definition to singular spaces.

\begin{definition}[label=\labelNFM{def:quasi-psh-function}, status=draft, tag=0000000]
    A \emph{quasi-plurisubharmonic function}, or \emph{quasi-psh function},
    is a function locally of the form $u = v + \rho$,
    where $v$ is psh and $\rho$ is smooth and real-valued.
    When viewing $u$ as a current, we assume that $u$ is not identically $-\infty$ on any connected component,
    so that $u \in L^1_{\mathrm{loc}}(X)$.

    Let $\omega$ be a smooth closed real $(1,1)$-form.
    We say that a quasi-psh function $u$ is \emph{$\omega$-psh} if
    \[
        \omega + \ddc u \geq 0
    \]
    as a current, and denote the set of these functions by $\operatorname{PSH}(X, \omega)$.
    Locally write $\omega = \ddc \rho$ with $\rho$ smooth and real-valued.
    Then $u$ is $\omega$-psh iff $u + \rho$ is psh.
\end{definition}

Under the hypotheses of \cref{prop:global-ddc-potentials}, assume $[T] = [\omega]$ and $T \geq 0$.
Then the global potential can be chosen $\omega$-psh.
If $T$ has continuous local potentials, this global potential is continuous.
Indeed, locally write $\omega = \ddc \rho$ with $\rho$ smooth and $T = \ddc v$ with $v$ psh.
Then $h \coloneqq u + \rho - v$ is a smooth pluriharmonic function,
so $u = v - \rho + h$ is quasi-psh and $\omega + \ddc u = T \geq 0$.
The local representatives agree on overlaps by uniqueness of psh representatives.
The same identity proves continuity when $v$ is continuous.

\begin{proposition}[label=\labelNFM{prop:wedge-with-psh-function}, status=draft, tag=0000000]
    Let $u$ be a locally bounded psh function on $X$ and $T$ a closed positive current of bidimension $(p,p)$,
    where $1 \leq p \leq d$.
    Since $T$ has measure coefficients and $u$ is a locally bounded Borel function,
    the current $uT$ is well-defined.
    Set
    \[
        \ddc u \wedge T \coloneqq \ddc(uT).
    \]
    Then $\ddc u \wedge T$ is a closed positive current of bidimension $(p - 1, p - 1)$.
    When $u$ is smooth, it agrees with the usual wedge product.
    See \cite[Chapter~III, \S 3, Proposition~3.2]{Demailly-2012-ComplexAnalyticDifferential}.
\end{proposition}
\begin{proof}
    The assertion is local.
    On a smaller coordinate ball, convolution gives smooth psh functions $u_j \downarrow u$.
    By local boundedness and dominated convergence for the coefficient measures of $T$,
    we have $u_j T \to uT$ weakly.
    Hence
    \[
        \ddc(uT) = \lim_{j \to \infty} \ddc(u_j T) = \lim_{j \to \infty} \ddc u_j \wedge T.
    \]
    Each current on the right is closed and positive, so the same holds for the limit.
\end{proof}

\begin{theorem}[label=\labelNFM{thm:bedford-taylor-products}, status=draft, tag=0000000]
    Let $T$ be a closed positive current of bidimension $(p,p)$,
    and let $u_1,\ldots,u_r$ be locally bounded psh functions with $1 \leq r \leq p \leq d$.
    The recursive formula
    \[
        \ddc u_1 \wedge \cdots \wedge \ddc u_r \wedge T \coloneqq \ddc \bigl(u_1 \, \ddc u_2 \wedge \cdots \wedge \ddc u_r \wedge T\bigr)
    \]
    defines a closed positive current of bidimension $(p - r, p - r)$,
    symmetric in $u_1,\ldots,u_r$.
    Moreover, if $u_{j,\ell} \downarrow u_j$ are decreasing sequences of psh functions for every $j$,
    then
    \[
        \ddc u_{1,\ell} \wedge \cdots \wedge \ddc u_{r,\ell} \wedge T \longrightarrow \ddc u_1 \wedge \cdots \wedge \ddc u_r \wedge T
    \]
    weakly.
    See \cite[Chapter~III, \S 3, Theorem~3.7 and Corollary~3.10]{Demailly-2012-ComplexAnalyticDifferential}.
\end{theorem}

\begin{construction}[label=\labelNFM{constr:products-of-currents-with-local-potentials}, status=draft, tag=0000000]
    Let $S_1,\ldots,S_r$ be closed positive currents of bidegree $(1,1)$ with locally bounded local potentials,
    and let $T$ be a closed positive current of bidimension $(p,p)$,
    where $1 \leq r \leq p \leq d$.
    Locally write $S_j = \ddc u_j$ with $u_j$ locally bounded and psh, and set
    \[
        S_1 \wedge \cdots \wedge S_r \wedge T \coloneqq \ddc u_1 \wedge \cdots \wedge \ddc u_r \wedge T.
    \]
    This is independent of the choice of local potentials.
    Indeed, two potentials differ by a smooth pluriharmonic function $h$,
    and $\ddc(hR) = 0$ for every closed current $R$ of pure bidimension.
    Together with symmetry in the preceding theorem, this proves independence in every factor.
    Hence the local products agree on overlaps and define a global closed positive current,
    symmetric in $S_1,\ldots,S_r$.

    In particular, for a closed positive current $S$ of bidegree $(1,1)$ with continuous local potentials,
    taking $T = 1$ gives the powers $S^r$ for $1 \leq r \leq d$,
    and $S^d$ is a positive measure.
    We will use this construction for Green currents and bifurcation currents below.
\end{construction}

We now allow the psh function to be unbounded.
The point is whether it is integrable against the coefficient measures of the current.

\begin{proposition}[label=\labelNFM{prop:wedge-with-integrable-psh-function}, status=draft, tag=0000000]
    Let $u \in \operatorname{PSH}(X) \cap L^1_{\mathrm{loc}}(X)$ and let $T$ be a closed positive current of bidimension $(p,p)$,
    where $1 \leq p \leq d$.
    In a coordinate chart, put $\beta \coloneqq \ddc |z|^2$ and let $\sigma_T \coloneqq T \wedge \beta^p$ be the trace measure of $T$.
    Suppose that $u \in L^1_{\mathrm{loc}}(\sigma_T)$ in every coordinate chart.
    Then $uT$ is a current with locally finite measure coefficients, and
    \[
        \ddc u \wedge T \coloneqq \ddc(uT)
    \]
    is a closed positive current of bidimension $(p - 1, p - 1)$.
\end{proposition}
\begin{proof}
    The coefficient measures of $T$ are locally dominated by a constant multiple of $\sigma_T$,
    so the integrability assumption defines $uT$.
    On a smaller coordinate ball, take smooth psh functions $u_j \downarrow u$.
    On each compact subset, $|u_j|$ is bounded by $|u| + C$ for some constant $C$.
    Dominated convergence gives $u_j T \to uT$ weakly.
    Hence
    \[
        \ddc(uT) = \lim_{j \to \infty} \ddc(u_j T) = \lim_{j \to \infty} \ddc u_j \wedge T,
    \]
    which is closed and positive as before.
\end{proof}

Here is a useful sufficient condition for this integrability.
Let $L(u)$ be the closed set of points where $u$ is not locally bounded.
It contains $\{u = -\infty\}$ and can be larger.
If $L(u) \cap \Supp T$ is discrete and $p \geq 1$,
then $u \in L^1_{\mathrm{loc}}(\sigma_T)$ and the preceding product is defined;
see \cite[Chapter~III, \S 4, Proposition~4.1]{Demailly-2012-ComplexAnalyticDifferential}.
For example, $u(z) = \log \|z\|$ on $\bbC^d$ has $L(u) = \{0\}$.
Applying this criterion recursively defines $(\ddc u)^r$ for $1 \leq r \leq d$.

\begin{remark}[label=\labelNFM{rmk:wedge-of-ddc-functions}, status=draft, tag=0000000]
    In general, the wedge product of two currents is not defined.
    For example, let $H = \{z_1 = 0\} \subset \bbC^2$ and consider the integration current $S = \delta_H$.
    By the Poincar\'e--Lelong formula in our normalization,
    \[
        S = \ddc \log |z_1|.
    \]
    But $\log |z_1| = -\infty$ on $H$, so $(\log |z_1|)S$ is not defined,
    and the preceding construction does not define $S \wedge S$.
    Here $u = \log |z_1|$ is locally integrable with respect to Lebesgue measure,
    but is not integrable against the trace measure of $S$, which is supported on $H$.
    Thus local integrability on $X$ alone does not suffice for forming $uT$.
    For the Poincar\'e--Lelong formula, see \cite[Chapter~III, \S 2, (2.15)]{Demailly-2012-ComplexAnalyticDifferential}.
\end{remark}



\subsection[status=draft,tag=0000000,label=\labelNFM{subsec:review-of-metrics}]{Hermitian line bundles}

Recall that we have the exponential sequence
\[ 0 \to \underline{\bbZ} \xrightarrow{n \mapsto 2n\pi\im } \scrO_X \xrightarrow{\exp} \scrO_X^* \to 0\]
and then the long exact sequence
\[ \cdots \to H^1(X,\scrO_X^*) \xrightarrow{c_1} H^2(X, \underline{\bbZ}) \to H^2(X,\scrO_X) \to \cdots \]

Given a line bundles $L$ on $X$, we can get a cohomological class $c_1(L) \in H^{1,1}(X,\bbR) \cap H^2(X,\bbZ)$.
\Yang{}
But it is not a differential forms.
We hope to find a canonical differential form $\alpha$ to represent it, i.e., $[\alpha] = c_1(L)$.
For this purpose, we need some more finer structure on the line bundle, namely, the \emph{metric}.

\begin{definition}[label=\labelNFM{def:hermitian-line-bundles}, status=draft, tag=0000000]
	Let $X$ be a complex analytic space and $E$ be a vector bundle on $X$.
	Fix a class of functions $\cinf \subset \scrF \subset Map()$, typically $\scrF=\scrC_{X,\bbC}^p$ or $\scrL_{X,\bbC}^p$.
	A \emph{$\scrF$-hermitian metric} on $E$ is a $\scrF$-ly varying family of hermitian inner products \(\langle \cdot, \cdot \rangle_x\) on the fibers \(E_x\) for each \(x \in X\),
	i.e.,
	\[\langle \cdot, \cdot \rangle_x: E_x \times E_x \to \mathbb{C}\]
	is a hermitian inner product for each \(x\), and for any local $\scrF$-sections \(s,t\) of \(E\), the function
	\[x \mapsto \langle s(x), t(x) \rangle_x\]
	is of class $\scrF$ on \(X\).
	A complex vector bundle equipped with a hermitian metric is called a \emph{hermitian vector bundle}.
	\Yang{}
\end{definition}

\begin{remark}\label{rmk:Riemannian_metric_induced_by_hermitian_metric}
	Let \(h\) be a smooth hermitian metric on a complex manifold \(X\).
	Then \(h\) induces a Riemannian metric \(g\) on the underlying real manifold of \(X\) by
	\[g(u,v) = \Real(h(u,v))\]
	for real tangent vectors \(u,v \in T_x X\).
\end{remark}

\begin{remark}[label=\labelNFM{rmk:metric-on-tangent-gives-a-form}, status=draft, tag=0000000]
	Let $h$ be a smooth hermitian metric on $TX$.
	Then it gives a smooth section in $\wedge^2 T^*X$ and hence a smooth $2$-forms in $\scrA^2_X(X)$.
	Moreover, it is a $(1,1)$ form.

	\Yang{}
\end{remark}


On this note, we will focus the case of hermitian line bundles.


%
% \begin{example}\label{eg:Fubini-Study_metric_on_P^n}
% 	Let \(\bbP^n\) be the complex projective space (parameterizing lines).
% 	Recall that \(\bbP^n(\bbC) \cong (\bbC^{n+1} \setminus \{0\})/\bbC^*\).
% 	We focus on the underlying smooth manifold structure.
% 	We have \(\bbC^* \cong S^1 \times \bbR_{>0}\) and \((\bbC^{n+1} \setminus \{0\})/ \bbR_{>0} \cong S^{2n+1}\).
% 	Hence \(\bbP^n(\bbC) \cong S^{2n+1}/S^1\).
% 	Note that \(S^1 \acts S^{2n+1} \subset \bbC^{n+1}\) is isometric with respect to the standard hermitian metric on \(\bbC^{n+1}\).
% 	Hence the quotient \(\bbP^n(\bbC)\) inherits a natural hermitian metric \(h_{\text{FS}}\), called the \emph{Fubini-Study metric}.
%
% 	On the standard affine chart \(U_i = \{[z_0: \cdots : z_n] \in \bbP^n \mid z_i \neq 0\}\) with coordinates \(z_{j,i} = z_j/z_i\) for \(j \neq i\),
% 	we know that \(T\bbP^n|_{U_i}\) is spanned by \(\{\partial_{j,i} = \partial/\partial z_{j,i}\}_{j \neq i}\).
% 	The Fubini-Study metric is given by
% 	\[h_{\text{FS}}(z_{-,i})(\partial_{k,i}, \partial_{l,i}) = \frac{\delta_{kl}}{1 + \sum_{r \neq i} |z_{r,i}|^2} - \frac{\overline{z_{k,i}} z_{l,i}}{(1 + \sum_{r \neq i} |z_{r,i}|^2)^2}.\]
%
% \Yang{relate this two FS metric by some sequence? like Euler sequence?}
% 	\Yang{}
% \end{example}

\begin{example}[label=\labelNFM{eg:hermitian-metric-on-trivial-bundle}, status=draft, tag=0000000]
	Let $X$ be a complex analytic variety and $L = \bbC \times X$ by the trivial line bundle.
	The standard hermitian inner product gives a  hermitian metric on $L$.
	Use the language of sheaves, we have
	\[ h_{tri}: \cinf_{X,\bbC} \times \cinf_{X,\bbC} \to \bbC, \quad (f,g) \mapsto f\bar{g}. \]
	More generally, given $\varphi \in \cinf_{X,\bbR}(X)$, we can construct a hermitian metric by
	\[ h_{\varphi}: \cinf_{X,\bbC} \times \cinf_{X,\bbC} \to \bbC, \quad (f,g) \mapsto \exp(\varphi) \cdot f\bar{g}. \]
\end{example}

\begin{example}\label{eg:Fubini-Study_metric_on_O_1}
	Let \(\bbP^n\) be the complex projective space in the sense of Grothendieck, i.e., it equips a universal quotient $\scrO^{\oplus n+1} \surjmap \scrO(1) \to 0$.
	Recall that on $U_k = \{X_k\neq 0\} \cong \bbA^n(\bbC)$ with coordinates $(x_i/x_k)_{i=0,\cdots,\hat{k},\cdots,n}$, $\scrO(1)$ is given by
	\[ \Gamma(U_k, \scrO(1); \cinf_{\bbP^n,\bbC}) = \Gamma(U_k, \scrO^{\oplus n+1}; \cinf_{\bbP^n,\bbC})/\sim \ = \left\{ \sum_{k} f_k X_k \;\middle|\; f_k \in \cinf_{X,\bbC}(U) \right\} \big/ \left(\frac{x_i}{x_k} \cdot X_k \sim X_i, \ i=0,\cdots,n \right). \]
	On $\scrO_X^{\oplus n+1} = \bigoplus_{i=0}^n \scrO_X \cdot X_i$, we have the standard hermitian metric
	\[ \left< \sum a_i X_i, \sum b_i X_i \right> \coloneqq \sum a_i \bar{b_i}, \qquad a_i,b_i \in \cinf_{X,\bbC}(U). \]
	Use the metric, the quotient $\scrO^{\oplus n+1} \xrightarrow{\pi} \scrO(1) \to 0$ $\cinf$-ly splits as
	\[ \scrO^{\oplus n+1} \cong \scrE \oplus \scrO(1), \qquad \scrE = \Ker \pi, \qquad \scrO(1) \cong (\scrE)^{\perp}\subset \scrO^{\oplus n+1} .\]
	The standard hermitian metric restricts to a hermitian metric $h_{FS}$ on $\scrO(1)$.
	This metric is called the \emph{Fubini-Study} metric.

	Let us focus on the affine piece $U = \{X_0\neq 0\} \cong \bbA^n(\bbC)$ with coordinates $(z_i = x_i/x_0)_{i=1,\cdots,n}$.
	As $\cinf_{\bbP^n,\bbC}(U)$-modules, $\Gamma(U, \scrE; \cinf_{\bbP^n,\bbC})$ has a basis $\{E_i = z_i \cdot X_0 - X_i\}_{i=0,\cdots,\hat{k},\cdots,n}$.
	Set $F = X_0 + \sum_{i=1}^n \overline{z_i} X_i \in \Gamma(U, \scrO^{\oplus n+1}; \cinf_{\bbP^n,\bbC})$.
	By some computation, we find that
	\[ \left<F, E_i \right> =  \left< z_i \cdot X_0 - X_i, X_0 + \sum_{i=1}^n \overline{z_i} X_i \right> = z_i - z_i = 0, \quad \forall i = 1,\cdots, n. \]
	And we have $\left<F,F\right> = 1 + \sum |z_i|^2$.
	Since $F|_x \neq 0 \in \scrO^{\oplus n+1}|_x$ for every $x \in U$, we have $\scrO(1)|_U \cong (\scrO \cdot F)|_U$.

	Under the basis $F,E_1,\cdots,E_n$, we have
	\[ X_0 = \frac{1}{1+\sum_{j=1}^n|z_j|^2} F + \sum_{i=1}^n \frac{\overline{z_i}}{1+\sum_{j=1}^n|z_j|^2} E_i = \frac{|x_0|^2}{|x_0|^2+\cdots+|x_n|^2} F + \sum f_i E_i. \]
	Hence
	\[ h_{FS}(X_0,X_0) = \left< \frac{|x_0|^2}{|x_0|^2+\cdots+|x_n|^2} F, \frac{|x_0|^2}{|x_0|^2+\cdots+|x_n|^2} F \right> = \frac{|x_0|^2}{|x_0|^2+\cdots+|x_n|^2}.\]
\end{example}


\begin{definition}[label=\labelNFM{def:tensor-product-of-metrics}, status=draft, tag=0000000]
	Let $X$ be a complex analytic variety and $\bar{L} = (L,\|\cdot\|)$, $\bar{L}_1 = (L_1,\|\cdot\|_1)$, $\bar{L}_2 = (L_2,\|\cdot\|_2)$ be metrized line bundles.

	On the dual $L^{-1}$, there is a natural metric given by $\|f\|^{\vee} = |f(s)|/\|s\|$, for any $s \neq 0$.

	On the tensor product $L_1 \otimes L_2$, there is a natural metric given by $s_1\otimes s_2 \mapsto \|s_1\|\|s_2\|$.
	With this metric, we define the \emph{tensor product of metrized line} as $\bar{L}_1 \otimes \bar{L}_2 = (L_1\otimes L_2, \|\cdot\|_1\otimes \|\cdot\|_2)$.

	Under these operators, the isomorphic class of line bundles forms a group, denoted by $\aPic(X)$.
\end{definition}


\begin{proposition}[label=\labelNFM{prop:hermitian-metric-on-an-complex-space}, status=draft, tag=0000000]
	Let $X$ be a complex analytic variety.
	We have an exact sequence of abelian groups
	\[ 0 \to \cinf_{X,\bbR}(X) \xrightarrow{\exp} \aPic(X) \to \Pic(X) \to 0. \]
\end{proposition}
\begin{proof}
	\Yang{}
\end{proof}

Under the correspondence of sections in a line bundle and Cartier divisors, the metric will becomes the \emph{green functions}.

\begin{definition}\label{def:green-functions-for-a-divisor}
	Let $X$ be a complex analytic variety.
	Let $D = \{(U_i,f_i)\}$ be a Caritier divisor on $X$.
	A \emph{green function of divisor $D$} is a (real-valued) function $g \in \scrC_{X,\bbR}(X \setminus \Supp D)$ such that $g|_{U_i} + \log |f_i|$ can be extended to a continous function on $U_i$.
	The pair $(D,g)$ is called a \emph{green divisor}.

	A green function is called \emph{prinicple} if it is of form $(\divisor f, -\log|f|)$ with $f \in \scrM_X^*$.

	Given two green divisor $(D_1,g_1)$ and $(D_2,g_2)$, we define $(D_1,g_1) + (D_2,g_2) \coloneqq (D_1+D_2,g_1+g_2)$ and $-(D_1,g_1) \coloneqq (-D_1,-g_1)$.
\end{definition}
\begin{remark}[label=\labelNFM{rmk:name-of-green-function}, status=draft, tag=0000000]
	Why such $g$ is called a green function? the Lelong equation.
	\Yang{}
\end{remark}

\begin{proposition}[label=\labelNFM{prop:group-of-green-divisors}, status=draft, tag=0000000]
	Let $X$ be a complex analytic variety.
	The green divisors from a groups, we have exact sequence


\end{proposition}


Let \(X\) be a complex analytic variety and $\calL$ a holomorphic line bundle on $X$.

\begin{theorem}[label=\labelNFM{thm:metrics-green-functions-and-forms}, status=draft, tag=0000000]
	Let \(X\) be a complex projective variety and  \Yang{}


	\begin{enumerate}
		\item The natural homomorphism, $\aCl(X) \to \aPic(X), (D,g) \mapsto (\scrO(D),\|\cdot\|_g)$, is an isomorphism.
		\item there is a chern class map $c_1: \aPic(X) \to \scrA^{1,1}_{\bbR}(X)$ with $c_1(L) = [c_1(L,\|\cdot\|)]$.
		\item we have the Lelong equation $\upd \upd^c g = \frac{\im}{\pi} \partial \overline{\partial} g = \omega_h - \delta_D$.
	\end{enumerate}

\end{theorem}

































